\documentclass[11pt]{article}
\usepackage{setspace}
\usepackage{geometry}
\geometry{margin=1in}
\setstretch{1.15}

\begin{document}

% --- Applicant Information ---
\noindent
\textbf{Zhongli Wang} \\
Phone: +49 160 9541 6015 \\
Email: nkwangzhongli@gmail.com \\
Website: https://sites.google.com/view/zhongli-wang \\
Bielefeld, Germany \\

\vspace{0.5cm}

% --- Recipient Information ---
\noindent
\textbf{TIAS Selection Committee} \\
Turku Institute for Advanced Studies \\
University of Turku \\
Turku, Finland \\

\vspace{0.5cm}

% --- Subject Line ---
\noindent
\textbf{Application for the TIES2 Senior Researcher Fellowship} \\

\vspace{0.3cm}

\noindent Dear Members of the TIAS Selection Committee,\\[0.2cm]

I am writing to apply for the TIES2 Fellowship at the Turku Institute for Advanced Studies. 
I am currently completing my Ph.D.\ in Economics at Bielefeld University and Università Cattolica del Sacro Cuore through the Marie Skłodowska–Curie EPOC program. 
My research examines AI adoption, network diffusion, and technology-mediated behavior, with a strong methodological orientation toward agent-based modeling, learning dynamics, and data-driven analysis using machine learning and large language models (LLMs). 
The TIES2 Fellowship offers an ideal environment to develop my proposed project, \emph{AI, Networks, and the Diffusion of Digital Technologies in Nordic Societies}, which aligns closely with TIAS’s mission to support independent, interdisciplinary, and societally relevant research.

Across my doctoral work, which I began in 2021, I have increasingly focused on how individuals and organizations learn, make decisions, and adapt in AI-rich environments.  
My dissertation develops computational models of heterogeneous agents who form beliefs, update strategies, and adapt behavior under information frictions and digital transformation.  
These models draw on network theory, diffusion dynamics, and behavioral economics to examine how technologies spread through populations and how structural conditions shape aggregate patterns.  
This experience has prepared me to investigate the Nordic context of AI diffusion, where network structures, institutional design, and digital infrastructures play a critical role in shaping household and firm behavior.

In addition to computational modeling, my research includes work on machine learning and human–AI interaction.  
One coauthored project, \emph{Polytope Fraud Theory}, uses ML to extract behavioral signals from large-scale financial data, demonstrating my ability to work with high-dimensional datasets, feature engineering, and interpretability tools.  
Another project develops an agentic workflow integrating LLMs into empirical research pipelines, providing insight into how AI systems reshape reasoning processes, knowledge acquisition, and user behavior.  
These experiences strengthen my ability to combine empirical, computational, and conceptual approaches—skills essential for the interdisciplinary agenda of the TIES2 fellowship.

The University of Turku offers a uniquely suitable environment for this work.  
The Economics and Management faculties host scholars working on digitalization, public policy, household behavior, and networked social systems, providing valuable connections for collaborative scientific exchange.  
My proposed secondment at Lucidminds AI (Netherlands) adds an applied, intersectoral dimension by enabling the development of transferable simulation tools, explainable AI components, and scenario-based policy prototypes.  
This partnership will ensure that the project contributes simultaneously to academic knowledge, societal relevance, and the design of evidence-informed digital strategies.

During the fellowship, I aim to contribute to the intellectual community at TIAS by engaging in interdisciplinary seminars, sharing methodological tools, and offering collaborative support to researchers working on AI, behavioral modeling, and digital society.  
I am particularly interested in fostering dialogue between economic modeling, network science, computational social science, and policy-oriented research—domains strongly represented at the University of Turku.

Moreover, the proposed project builds directly on the methodological foundations I have developed throughout my doctoral work. 
I have extensive experience implementing agent-based and network diffusion models, estimating behavioral parameters using simulation-based econometrics, and integrating machine-learning tools—including large language models—into empirical and computational workflows. 
These skills position me to construct diffusion frameworks that incorporate heterogeneous households and firms, social influence, institutional frictions, and the digital infrastructures characteristic of Nordic societies. 
My prior work and interests with large-scale and multimodal datasets (such as geospatial information, digital behavioral traces, and high-dimensional text corpora) equips me to analyze adoption patterns using both statistical and ML-based approaches. 
Because all chapters of my dissertation are single-authored, I am accustomed to independently designing complex model architectures, validating behavioral mechanisms, and translating theoretical constructs into reproducible, well-documented code. 
Together, these experiences ensure that I have both the conceptual and technical capacity to execute the proposed project with the rigor, originality, and interdisciplinary depth expected of a TIAS Fellow.

Thank you so much for considering my application! 
I would welcome the opportunity to discuss how my research agenda, methodological background, and intersectoral experience can contribute to TIAS and to the University of Turku’s broader strengths in digital society research, interdisciplinary collaboration, and open science.

\vspace{6mm}

\noindent Sincerely,\\[15mm]
\noindent Zhongli Wang

\end{document}
